%!TEX root = forallxyyc.tex
\chapter[Referência rápida]{Referência rápida}
%\pagestyle{plain}
\section{Tabelas-verdade características}
\label{app.CharacteristicTTs}

\begin{tabular}{c|c}
\metav{A} & \enot\metav{A}\\
\hline
T & F\\
F & T \\
\phantom{.}\\
\phantom{.}
\end{tabular}
\hfill
\begin{tabular}{c c|c|c|c|c}
\metav{A} & \metav{B} & $\metav{A}\eand\metav{B}$ & $\metav{A}\eor\metav{B}$ & $\metav{A}\eif\metav{B}$ & $\metav{A}\eiff\metav{B}$\\
\hline
T & T & T & T & T & T\\
T & F & F & T & F & F\\
F & T & F & T & T & F\\
F & F & F & F & T & T
\end{tabular}


\vfill

\section{Simbolização}
\begin{center}
\label{app.symbolization}
\begin{tabular*}{\textwidth}{rl}
\multicolumn{2}{c}{\textsc{Conectivos sentenciais}}\\ \\
Não é o caso de que $P$ & $\enot P$\\
$P$ ou $Q$ & $(P \eor Q)$\\
Nem $P$ nem $Q$ & $\enot(P \eor Q)$\ ou \ $(\enot P \eand \enot Q)$\\
$P$ e $Q$ & $(P \eand Q)$\\
Se $P$, então $Q$ & $(P \eif Q)$\\
$P$ somente se $Q$ & $(P \eif Q)$\\
$P$ se, e somente se, $Q$ & $(P \eiff Q)$\\
$P$ a menos que $Q$ & $(\enot Q \eif P)$ or $(P \eor Q)$\\[12pt]
\multicolumn{2}{c}{\label{SymbolizingPredicates}\textsc{Predicados}}\\ \\
Todos os $F$s são $G$s & $\forall x(\atom{F}{x} \eif \atom{G}{x})$\\
Alguns $F$s são $G$s & $\exists x(\atom{F}{x} \eand \atom{G}{x})$\\
Nem todos os $F$s são $G$s & $\enot\forall x(\atom{F}{x} \eif \atom{G}{x})$\ or\\
& $\exists x(\atom{F}{x} \eand \enot \atom{G}{x})$\\
Nenhum $F$ é $G$ & $\forall x(\atom{F}{x} \eif\enot \atom{G}{x})$\ or\\
& $\enot\exists x(\atom{F}{x} \eand \atom{G}{x})$\\
Somente $F$s são $G$s & $\forall x(\atom{G}{x} \eif \atom{F}{x})$\\
& $\enot\exists x(\enot\atom{F}{x} \eand \atom{G}{x})$\\[12pt]
\multicolumn{2}{c}{\textsc{Identidade}}\\ \\
Somente $c$ é $G$ & $\forall x(\atom{G}{x} \eiff x=c)$\\
Tudo o que não é &\\ que $c$
é $G$ & $\forall x(\enot\, {x = c} \eif \atom{G}{x} )$\\
Tudo, exceto $c$
 é $G$ & $\forall x(\enot\, {x = c} \eiff \atom{G}{x} )$\\
%$j$ é mais $R$ do que qualquer outra coisa. & $\forall x(x\neq j \eif Rjx)$\\
O $F$ é $G$ & $\exists x(\atom{F}{x} \eand \forall y(\atom{F}{y} \eif x=y) \eand \atom{G}{x} )$\\
Não é o caso de que
 o $F$ é $G$ & $\enot\exists x(\atom{F}{x} \eand \forall y(\atom{F}{y} \eif x=y) \eand \atom{G}{x} )$\\
O $F$ não é $G$ & $\exists x(\atom{F}{x} \eand \forall y(\atom{F}{y} \eif x=y) \eand \enot \atom{G}{x} )$
\end{tabular*}
\end{center}






% BEGIN: symbolizing cardinality

\newpage
\section{Usando a identidade para simbolizar quantidades}

\subsection*{Há pelo menos \blank\ $F$s.}
\label{summary.atleast}


\ifHTMLtarget
\begin{tabular*}{\textwidth}{rl}
	um & $\exists x\,\atom{F}{x}$\\
	dois & $\exists x_1\exists x_2(\atom{F}{x_1} \eand \atom{F}{x_2}
	\eand \enot x_1  = x_2)$\\
	três & $\exists x_1\exists x_2\exists x_3(\atom{F}{x_1} \eand \atom{F}{x_2} \eand \atom{F}{x_3} \eand \enot x_1 = x_2 \eand\enot x_1 = x_3 \eand \enot x_2 = x_3)$\\
	quatro & $\exists x_1\exists x_2\exists x_3\exists x_4 (\atom{F}{x_1} \eand \atom{F}{x_2} \eand \atom{F}{x_3} \eand \atom{F}{x_4} \eand \enot x_1 = x_2 \eand \enot x_1 = x_3 \eand \enot x_1 = x_4 \eand \enot x_2 = x_3 \eand \enot x_2 = x_4 \eand \enot x_3 = x_4)$\\
	$n$ & $\exists x_1\ldots\exists x_n(\atom{F}{x_1} \eand \ldots \eand \atom{F}{x_n} \eand \enot x_1 = x_2 \eand\ldots\eand \enot x_{n-1} = x_n)$
\end{tabular*}
\else
\begin{tabular*}{\textwidth}{rl}
	um & $\exists x\,\atom{F}{x}$\\
	dois & $\exists x_1\exists x_2(\atom{F}{x_1} \eand \atom{F}{x_2}
	\eand \enot x_1  = x_2)$\\
	três & $\exists x_1\exists x_2\exists x_3(\atom{F}{x_1} \eand \atom{F}{x_2} \eand \atom{F}{x_3} \eand {}$\\
	& \quad$\enot x_1 = x_2 \eand\enot x_1 = x_3 \eand \enot x_2 = x_3)$\\
	quatro & $\exists x_1\exists x_2\exists x_3\exists x_4 (\atom{F}{x_1} \eand \atom{F}{x_2} \eand \atom{F}{x_3} \eand \atom{F}{x_4} \eand {}$\\
	& \quad $\enot x_1 = x_2 \eand \enot x_1 = x_3 \eand \enot x_1 = x_4 \eand {}$\\
	& \qquad$ \enot x_2 = x_3 \eand \enot x_2 = x_4 \eand \enot x_3 = x_4)$\\
	$n$ & $\exists x_1\ldots\exists x_n(\atom{F}{x_1} \eand \ldots \eand \atom{F}{x_n} \eand {}$\\
	& \quad$\enot x_1 = x_2 \eand\ldots\eand \enot x_{n-1} = x_n)$
\end{tabular*}
\fi

\subsection*{Há no máximo \blank\ $F$s.}
\label{summary.atmost}

Uma maneira de dizer `há no máximo $n$ $F$s' é colocar um sinal de negação antes da simbolização de `há pelo menos $n+1$ $F$s'. Equivalentemente, podemos oferecer:

\noindent
\ifHTMLtarget
\begin{tabular*}{\textwidth}{rl}
	um & $\forall x_1\forall x_2\bigl[(\atom{F}{x_1} \eand
	\atom{F}{x_2}) \eif x_1=x_2\bigr]$\\
	dois & $\forall x_1\forall x_2\forall x_3\bigl[(\atom{F}{x_1} \eand \atom{F}{x_2} \eand \atom{F}{x_3}) \eif (x_1=x_2 \eor x_1=x_3 \eor x_2=x_3)\bigr]$\\
	três & $\forall x_1\forall x_2\forall x_3\forall x_4\bigl[(\atom{F}{x_1} \eand \atom{F}{x_2} \eand \atom{F}{x_3} \eand \atom{F}{x_4}) \eif (x_1=x_2 \eor x_1=x_3 \eor x_1=x_4 \eor x_2=x_3 \eor x_2=x_4 \eor x_3=x_4)\bigr]$\\
	$n$ & $\forall x_1\ldots\forall x_{n+1} \bigl[(\atom{F}{x_1} \eand \ldots \eand \atom{F}{x_{n+1}}) \eif (x_1=x_2 \eor \ldots \eor x_n=x_{n+1})\bigr]$
\end{tabular*}
\else
\begin{tabular*}{\textwidth}{rl}
	um & $\forall x_1\forall x_2\bigl[(\atom{F}{x_1} \eand
	\atom{F}{x_2}) \eif x_1=x_2\bigr]$\\
	dois & $\forall x_1\forall x_2\forall x_3\bigl[(\atom{F}{x_1} \eand \atom{F}{x_2} \eand \atom{F}{x_3}) \eif {}$\\ & \quad$(x_1=x_2 \eor x_1=x_3 \eor x_2=x_3)\bigr]$\\
	três & $\forall x_1\forall x_2\forall x_3\forall x_4\bigl[(\atom{F}{x_1} \eand \atom{F}{x_2} \eand \atom{F}{x_3} \eand \atom{F}{x_4}) \eif {}$\\
	& \quad$(x_1=x_2 \eor x_1=x_3 \eor x_1=x_4 \eor {}$\\
	& \qquad$x_2=x_3 \eor x_2=x_4 \eor x_3=x_4)\bigr]$\\
	$n$ & $\forall x_1\ldots\forall x_{n+1}
	\bigl[(\atom{F}{x_1} \eand \ldots \eand \atom{F}{x_{n+1}}) \eif {}$\\
	& \quad$(x_1=x_2 \eor \ldots \eor x_n=x_{n+1})\bigr]$
\end{tabular*}
\fi


\subsection*{Há exatamente \blank\ $F$s.}
\label{summary.exactly}

Uma maneira de dizer `há exatamente $n$ $F$s' é conjungir duas das simbolizações acima e dizer `há pelo menos $n$ $F$s e há no máximo $n$ $F$s'. As fórmulas equivalentes a seguir são mais curtas:

\noindent
\ifHTMLtarget
\begin{tabular*}{\textwidth}{rl}
	zero & $\forall x\,\enot \atom{F}{x}$\\
	um & $\exists x\bigl[\atom{F}{x} \eand \forall y(\atom{F}{y} \eif x = y)\bigr]$\\
	dois & $\exists x_1\exists x_2\bigl[\atom{F}{x_1} \eand \atom{F}{x_2} \eand \enot x_1 = x_2 \eand \forall y\bigl(\atom{F}{y} \eif (y= x_1 \eor y = x_2)\bigr) \bigr]$\\
	três & $\exists x_1\exists x_2\exists x_3\bigl[\atom{F}{x_1} \eand \atom{F}{x_2} \eand \atom{F}{x_3} \eand \enot x_1 =  x_2 \eand \enot  x_1 = x_3 \eand \enot x_2 = x_3 \eand \forall y\bigl(\atom{F}{y} \eif (y = x_1 \eor y = x_2 \eor y =  x_3)\bigr) \bigr]$\\
	$n$ & $\exists x_1\ldots\exists x_n\bigl[\atom{F}{x_1} \eand\ldots\eand \atom{F}{x_n}  \eand \enot x_1 = x_2 \eand\ldots\eand \enot x_{n-1}= x_n \eand \forall y\bigl(\atom{F}{y} \eif (y= x_1 \eor \ldots \eor y= x_n)\bigr)\bigr]$
\end{tabular*}
\else
\begin{tabular*}{\textwidth}{rl}
	zero & $\forall x\,\enot \atom{F}{x}$\\
	um & $\exists x\bigl[\atom{F}{x} \eand \forall y(\atom{F}{y} \eif x = y)\bigr]$\\
	dois & $\exists x_1\exists x_2\bigl[\atom{F}{x_1} \eand \atom{F}{x_2} \eand {}$\\
	& \quad$\enot x_1 = x_2 \eand \forall y\bigl(\atom{F}{y} \eif (y= x_1 \eor y = x_2)\bigr) \bigr]$\\
	três & $\exists x_1\exists x_2\exists x_3\bigl[\atom{F}{x_1} \eand \atom{F}{x_2} \eand \atom{F}{x_3} \eand {}$\\
	& \quad$\enot x_1 =  x_2 \eand \enot  x_1 = x_3 \eand \enot x_2 = x_3 \eand {}$\\
	& \qquad$\forall y\bigl(\atom{F}{y} \eif (y = x_1 \eor y = x_2 \eor y =  x_3)\bigr) \bigr]$\\
	$n$ & $\exists x_1\ldots\exists x_n\bigl[\atom{F}{x_1} \eand\ldots\eand \atom{F}{x_n}  \eand {}$\\
	& \quad$\enot x_1 = x_2 \eand\ldots\eand \enot x_{n-1}= x_n \eand {}$\\
	& \qquad$\forall y\bigl(\atom{F}{y} \eif (y= x_1 \eor \ldots \eor y= x_n)\bigr)\bigr]$
\end{tabular*}
\fi
%\item[one] $\exists x\forall y\bigl[\atom{F}{x} \eand (\atom{F}{y} \eif y = x)\bigr]$
%\item[two] $\exists x\exists y\forall z\Bigl(\atom{F}{x} \eand \atom{F}{y} \eand \bigl[\atom{F}{z} \eif (z=x \eor z=y)\bigr] \eand x \neq y\Bigr)$
%\item[three] $\exists x_1\exists x_2\exists x_3\forall y\Bigl(\atom{F}{x_1} \eand \atom{F}{x_2} \eand \atom{F}{x_3} \eand [\atom{F}{y} \eif (y=x_1 \eor y=x_2 \eor y=x_3)] \eand x_1 \neq x_2 \eand x_1 \neq x_3 \eand x_2 \neq x_3\Bigr)$
%\item[n] $\exists x_1\cdots\exists x_n\forall y\Bigl(\atom{F}{x_1} \eand \cdots \eand \atom{F}{x_n} \eand \bigl[\atom{F}{y} \eif (y=x_1 \eor \cdots \eor y=x_n)\bigr] \eand x_1 \neq x_2 \eand\cdots\eand x_{n-1}\neq x_n\Bigr)$


\label{ProofRules}
\newpage\section{Regras básicas de dedução para LVF}

% \mathem é uma caixa vazia da mesma largura que $m$; usada para alinhar as regras
% properly
\newlength{\mathemwd}
\settowidth{\mathemwd}{\mbox{\small$m$}}
\def\mathem{{\hbox to\mathemwd{}}}

% If compiling to PDF, make all the rules small(er) and avoid left margin
\ifHTMLtarget\else
% \renewenvironment{fitchproof}
% 	{\noindent\par\noindent\small$\begin{nd}}
% 	{\end{nd}$\noindent\normalsize\ignorespacesafterend}
%
\renewenvironment{fitchproof}[1][]
  {\begin{flushleft}\small\begin{nd}[#1]}
  {\end{nd}\end{flushleft}}
\fi
%{\LARGE \textbf{Regras básicas de dedução}}
\begin{multicols}{2}
\subsection*{Reiteração}

\begin{fitchproof}
	\have[m]{a}{\metav{A}}
	\have[\ ]{c}{\metav{A}} \by{R}{a}
\end{fitchproof}

\subsection*{Conjunção}

\begin{fitchproof}
	\have[m]{a}{\metav{A}}
	\have[n]{b}{\metav{B}}
	\have[\ ]{c}{\metav{A}\eand\metav{B}} \ai{a, b}
\end{fitchproof}
\begin{fitchproof}
	\have[m]{ab}{\metav{A}\eand\metav{B}}
	\have[\ ]{a}{\metav{A}} \ae{ab}
\end{fitchproof}
\begin{fitchproof}
	\have[m]{ab}{\metav{A}\eand\metav{B}}
	\have[\ ]{b}{\metav{B}} \ae{ab}
\end{fitchproof}

\subsection*{Condicional}

\begin{fitchproof}
	\open
		\hypo[m]{a}{\metav{A}}\AS
		\have[n]{b}{\metav{B}}
	\close
	\have[\ ]{ab}{\metav{A}\eif\metav{B}}\ci{a-b}
\end{fitchproof}
\begin{fitchproof}
	\have[m]{ab}{\metav{A}\eif\metav{B}}
	\have[n]{a}{\metav{A}}
	\have[\ ]{b}{\metav{B}} \ce{ab,a}
\end{fitchproof}

\subsection*{Negação}

\begin{fitchproof}
\open
	\hypo[m]{a}{\metav{A}}\AS
	\have[n]{nb}{\ered}
\close
\have[\ ]{na}{\enot\metav{A}}\ni{a-nb}
\end{fitchproof}
\begin{fitchproof}
	\have[m]{na}{\enot\metav{A}}
	\have[n]{a}{\metav{A}}
	\have[ ]{bot}{\ered}\ri{na, a}
\end{fitchproof}

\subsection*{Prova indireta}

\begin{fitchproof}
\open
	\hypo[i]{a}{\enot\metav{A}}\AS
	\have[j]{nb}{\ered}
\close
\have[\mathem]{na}{\metav{A}}\ip{a-nb}
\end{fitchproof}


\subsection*{Explosão}

\begin{fitchproof}
\have[m]{bot}{\ered}
\have[ ]{}{\metav{A}}\re{bot}
\end{fitchproof}

\subsection*{Disjunção}

\begin{fitchproof}
	\have[m]{a}{\metav{A}}
	\have[\ ]{ab}{\metav{A}\eor\metav{B}}\oi{a}
\end{fitchproof}
\begin{fitchproof}
	\have[m]{a}{\metav{A}}
	\have[\ ]{ba}{\metav{B}\eor\metav{A}}\oi{a}
\end{fitchproof}
\begin{fitchproof}
	\have[m]{ab}{\metav{A}\eor\metav{B}}
	\open
		\hypo[i]{a}{\metav{A}}\AS
		\have[j]{c1}{\metav{C}}
	\close
	\open
		\hypo[k]{b}{\metav{B}}\AS
		\have[l]{c2}{\metav{C}}
	\close
	\have[\mathem]{c}{\metav{C}} \oe{ab,a-c1, b-c2}
\end{fitchproof}

\subsection*{Bicondicional}

\begin{fitchproof}
	\open
		\hypo[i]{a1}{\metav{A}} \AS
		\have[j]{b1}{\metav{B}}
	\close
	\open
		\hypo[k]{b2}{\metav{B}}\AS
		\have[l]{a2}{\metav{A}}
	\close
	\have[\mathem]{ab}{\metav{A}\eiff\metav{B}}\bi{a1-b1,b2-a2}
\end{fitchproof}
\begin{fitchproof}
	\have[m]{ab}{\metav{A}\eiff\metav{B}}
	\have[n]{a}{\metav{A}}
	\have[\ ]{b}{\metav{B}} \be{ab,a}
\end{fitchproof}
\begin{fitchproof}
	\have[m]{ab}{\metav{A}\eiff\metav{B}}
	\have[n]{a}{\metav{B}}
	\have[\ ]{b}{\metav{A}} \be{ab,a}
\end{fitchproof}

\end{multicols}

\section{Regras derivadas para LVF}
\begin{multicols}{2}
\subsection*{Silogismo disjuntivo}
\begin{fitchproof}
	\have[m]{ab}{\metav{A} \eor \metav{B}}
	\have[n]{nb}{\enot \metav{A}}
	\have[\ ]{con}{\metav{B}}\by{DS}{ab, nb}
\end{fitchproof}
\begin{fitchproof}
	\have[m]{ab}{\metav{A} \eor \metav{B}}
	\have[n]{nb}{\enot \metav{B}}
	\have[\ ]{con}{\metav{A}}\by{DS}{ab, nb}
\end{fitchproof}

\subsection*{Modus Tollens}

\begin{fitchproof}
	\have[m]{ab}{\metav{A}\eif\metav{B}}
	\have[n]{a}{\enot\metav{B}}
	\have[\ ]{b}{\enot\metav{A}} \by{MT}{ab,a}
\end{fitchproof}

\subsection*{Eliminação da dupla negação}
	\begin{fitchproof}
		\have[m]{dna}{\enot \enot \metav{A}}
		\have[ ]{a}{\metav{A}}\dne{dna}
	\end{fitchproof}


\subsection*{Terceiro excluído}
	\begin{fitchproof}
		\open
			\hypo[i]{a}{\metav{A}}\AS
			\have[j]{c1}{\metav{B}}
		\close
		\open
			\hypo[k]{b}{\enot\metav{A}}\AS
			\have[l]{c2}{\metav{B}}
		\close
		\have[\mathem]{ab}{\metav{B}}\tnd{a-c1,b-c2}
	\end{fitchproof}

%
%\subsection*{Hypothetical Syllogism}
%
%\begin{fitchproof}
%	\have[m]{ab}{\metav{A}\eif\metav{B}}
%	\have[n]{bc}{\metav{B}\eif\metav{C}}
%	\have[\ ]{ac}{\metav{A}\eif\metav{C}}\by{HS}{ab,bc}
%\end{fitchproof}

\subsection*{Regras de De Morgan}
\begin{fitchproof}
	\have[m]{ab}{\enot (\metav{A} \eor \metav{B})}
	\have[\ ]{dm}{\enot \metav{A} \eand \enot \metav{B}}\dem{ab}
\end{fitchproof}
\begin{fitchproof}
	\have[m]{ab}{\enot \metav{A} \eand \enot \metav{B}}
	\have[\ ]{dm}{\enot (\metav{A} \eor \metav{B})}\dem{ab}
\end{fitchproof}
\begin{fitchproof}
	\have[m]{ab}{\enot (\metav{A} \eand \metav{B})}
	\have[\ ]{dm}{\enot \metav{A} \eor \enot \metav{B}}\dem{ab}
\end{fitchproof}
\begin{fitchproof}
	\have[m]{ab}{\enot \metav{A} \eor \enot \metav{B}}
	\have[\ ]{dm}{\enot (\metav{A} \eand \metav{B})}\dem{ab}
\end{fitchproof}
\end{multicols}

%\newpage

\section{Regras básicas de dedução para LPO}

\begin{multicols}{2}
\subsection*{Eliminação universal}

\begin{fitchproof}
	\have[m]{a}{\forall \metav{x}\,\metav{A}(\ldots \metav{x} \ldots \metav{x}\ldots)}
	\have[\ ]{c}{\metav{A}(\ldots \metav{c} \ldots \metav{c}\ldots)} \Ae{a}
\end{fitchproof}

\subsection*{Introdução universal}

\begin{fitchproof}
	\have[m]{a}{\metav{A}(\ldots \metav{c} \ldots \metav{c}\ldots)}
	\have[\ ]{c}{\forall \metav{x}\,\metav{A}(\ldots \metav{x} \ldots \metav{x}\ldots)} \Ai{a}
\end{fitchproof}
\begin{raggedright}
\metav{c} não pode ocorrer em nenhuma suposição não descarregada
\end{raggedright}
\end{multicols}

\begin{multicols}{2}
\subsection*{Introdução existencial}

\begin{fitchproof}
	\have[m]{a}{\metav{A}(\ldots \metav{c} \ldots \metav{c}\ldots)}
	\have[\ ]{c}{\exists \metav{x}\,\metav{A}(\ldots \metav{x} \ldots \metav{c}\ldots)}\Ei{a}
\end{fitchproof}

\vfill\null\columnbreak

\subsection*{Eliminação existencial}

\begin{fitchproof}
	\have[m]{a}{\exists \metav{x}\,\metav{A}(\ldots \metav{x} \ldots \metav{x}\ldots)}
	\open
		\hypo[i]{b}{\metav{A}(\ldots \metav{c} \ldots \metav{c}\ldots)}\AS
		\have[j]{c}{\metav{B}}
	\close
	\have[\ ]{d}{\metav{B}}\Ee{a,b-c}
\end{fitchproof}
\begin{raggedright}
\noindent \metav{c} não pode ocorrer em nenhuma suposição não descarregada, em $\exists \metav{x}\,\metav{A}(\ldots \metav{x} \ldots \metav{x}\ldots)$, ou em \metav{B}\end{raggedright}
\end{multicols}

\subsection*{Introdução da identidade}

\begin{fitchproof}
	\have[\mathem]{x}{\metav{c}=\metav{c}} \by{=I}{}
\end{fitchproof}


\subsection*{Eliminação da identidade}

\begin{multicols}{2}
\begin{fitchproof}
	\have[m]{e}{\metav{a}=\metav{b}}
	\have[n]{a}{\metav{A}(\ldots \metav{a} \ldots \metav{a}\ldots)}
	\have[\ ]{ea1}{\metav{A}(\ldots \metav{b} \ldots \metav{a}\ldots)} \by{=E}{e,a}
\end{fitchproof}
\begin{fitchproof}
	\have[m]{e}{\metav{a}=\metav{b}}
	\have[n]{a}{\metav{A}(\ldots \metav{b} \ldots \metav{b}\ldots)}
	\have[\ ]{ea2}{\metav{A}(\ldots \metav{a} \ldots \metav{b}\ldots)} \by{=E}{e,a}
\end{fitchproof}
\end{multicols}

%\begin{minipage}{\textwidth} % hack to keep section header with table
\section{Regras derivadas para LPO}

\begin{multicols}{2}
\begin{fitchproof}
	\have[m]{ab}{\forall \metav{x}\enot \metav{A}}
	\have[\ ]{ac}{\enot \exists \metav{x} \metav{A}}\cq{m}
\end{fitchproof}
\begin{fitchproof}
	\have[m]{ab}{\enot \exists \metav{x}  \metav{A}}
	\have[\ ]{ac}{\forall \metav{x}\enot\metav{A}}\cq{m}
\end{fitchproof}
\begin{fitchproof}
	\have[m]{ab}{\exists \metav{x}\enot\metav{A}}
	\have[\ ]{ac}{\enot \forall \metav{x} \metav{A}}\cq{m}
\end{fitchproof}
\begin{fitchproof}
	\have[m]{ab}{\enot \forall \metav{x}  \metav{A}}
	\have[\ ]{ac}{\exists \metav{x}\enot \metav{A}}\cq{m}
\end{fitchproof}
\end{multicols}
%\end{minipage}

\section{Regras para o uso de subprovas e citações}

Para citar uma linha individual ao aplicar uma regra:
\begin{compactlist}
\item a linha deve vir antes da linha em que a regra é aplicada, mas
\item não pode ocorrer dentro de uma subprova que tenha sido fechada antes da linha em que a regra é aplicada.
\end{compactlist}
Para citar uma subprova ao aplicar uma regra:
\begin{compactlist}
\item a subprova citada deve estar inteiramente antes da aplicação da regra em que é citada,
\item a subprova citada não pode estar dentro de outra subprova fechada que seja encerrada na linha em que é citada, e
\item a última linha da subprova citada não pode ocorrer dentro de uma subprova aninhada.
\end{compactlist}

\section{Regras para cadeias de equivalências}

\ifHTMLtarget
\begin{align*}
\lnot\lnot \metav{P} & \Leftrightarrow \metav{P} && \text{(DN)}\\[2ex]
(\metav{P} \eif \metav{Q}) & \Leftrightarrow (\lnot \metav{P} \lor \metav{Q})
&& \text{(Cond)}\\
\lnot(\metav{P} \eif \metav{Q}) & \Leftrightarrow (\metav{P} \land \lnot\metav{Q}) \\[2ex]
(\metav{P} \eiff \metav{Q}) & \Leftrightarrow ((\metav{P} \eif \metav{Q}) \land  (\metav{Q} \eif \metav{P}))
&& \text{(Bicond)}\\[2ex]
\lnot(\metav{P} \land \metav{Q}) & \Leftrightarrow (\lnot\metav{P} \lor \lnot\metav{Q})
&& \text{(DeM)}\\
\lnot(\metav{P} \lor \metav{Q}) & \Leftrightarrow (\lnot\metav{P} \land \lnot\metav{Q}) \\[2ex]
(\metav{P} \lor \metav{Q}) & \Leftrightarrow (\metav{Q} \lor \metav{P}) & &\text{(Comm)}\\
(\metav{P} \land \metav{Q}) & \Leftrightarrow (\metav{Q} \land \metav{P})\\[2ex]
(\metav{P} \land (\metav{Q} \lor \metav{R})) & \Leftrightarrow ((\metav{P} \land \metav{Q}) \lor (\metav{P} \land \metav{R}))
&& \text{(Dist)}\\
(\metav{P} \lor (\metav{Q} \land \metav{R})) & \Leftrightarrow ((\metav{P} \lor \metav{Q}) \land (\metav{P} \lor \metav{R}))\\
(\metav{P} \lor (\metav{Q} \lor \metav{R})) & \Leftrightarrow ((\metav{P} \lor \metav{Q}) \lor \metav{R}) & &\text{(Assoc)}\\
(\metav{P} \land (\metav{Q} \land \metav{R})) & \Leftrightarrow ((\metav{P} \land \metav{Q}) \land \metav{R})\\[2ex]
(\metav{P} \lor \metav{P}) & \Leftrightarrow \metav{P} && \text{(Id)}\\
(\metav{P} \land \metav{P}) & \Leftrightarrow \metav{P}\\[2ex]
(\metav{P} \land (\metav{P} \lor \metav{Q})) & \Leftrightarrow \metav{P}&& \text{(Abs)}\\
(\metav{P} \lor (\metav{P} \land \metav{Q})) & \Leftrightarrow \metav{P}\\[2ex]
(\metav{P} \land (\metav{Q} \lor \lnot\metav{Q})) & \Leftrightarrow \metav{P}  && \text{(Simp)}\\
(\metav{P} \lor (\metav{Q} \land \lnot\metav{Q})) & \Leftrightarrow \metav{P}\\[2ex]
(\metav{P} \lor (\metav{Q} \lor \lnot\metav{Q})) & \Leftrightarrow (\metav{Q} \lor \lnot\metav{Q}) \\
(\metav{P} \land (\metav{Q} \land \lnot\metav{Q})) & \Leftrightarrow (\metav{Q} \land \lnot\metav{Q})
\end{align*}
\else

\begin{align*}
\lnot\lnot \metav{P} & \Leftrightarrow \metav{P} && \text{(DN)}\\[2ex]
(\metav{P} \eif \metav{Q}) & \Leftrightarrow (\lnot \metav{P} \lor \metav{Q})
&& \text{(Cond)}\\
\lnot(\metav{P} \eif \metav{Q}) & \Leftrightarrow (\metav{P} \land \lnot\metav{Q}) \\[2ex]
(\metav{P} \eiff \metav{Q}) & \Leftrightarrow ((\metav{P} \eif \metav{Q}) \land  (\metav{Q} \eif \metav{P}))
&& \text{(Bicond)}\\[2ex]
\lnot(\metav{P} \land \metav{Q}) & \Leftrightarrow (\lnot\metav{P} \lor \lnot\metav{Q})
&& \text{(DeM)}\\
\lnot(\metav{P} \lor \metav{Q}) & \Leftrightarrow (\lnot\metav{P} \land \lnot\metav{Q}) \\[2ex]
(\metav{P} \lor \metav{Q}) & \Leftrightarrow (\metav{Q} \lor \metav{P}) & &\text{(Comm)}\\
(\metav{P} \land \metav{Q}) & \Leftrightarrow (\metav{Q} \land \metav{P})\\[2ex]
(\metav{P} \land (\metav{Q} \lor \metav{R})) & \Leftrightarrow ((\metav{P} \land \metav{Q}) \lor (\metav{P} \land \metav{R}))
&& \text{(Dist)}\\
(\metav{P} \lor (\metav{Q} \land \metav{R})) & \Leftrightarrow ((\metav{P} \lor \metav{Q}) \land (\metav{P} \lor \metav{R}))\\
(\metav{P} \lor (\metav{Q} \lor \metav{R})) & \Leftrightarrow ((\metav{P} \lor \metav{Q}) \lor \metav{R}) & &\text{(Assoc)}\\
(\metav{P} \land (\metav{Q} \land \metav{R})) & \Leftrightarrow ((\metav{P} \land \metav{Q}) \land \metav{R})\\[2ex]
(\metav{P} \lor \metav{P}) & \Leftrightarrow \metav{P} && \text{(Id)}\\
(\metav{P} \land \metav{P}) & \Leftrightarrow \metav{P}\\[2ex]
(\metav{P} \land (\metav{P} \lor \metav{Q})) & \Leftrightarrow \metav{P}&& \text{(Abs)}\\
(\metav{P} \lor (\metav{P} \land \metav{Q})) & \Leftrightarrow \metav{P}\\[2ex]
(\metav{P} \land (\metav{Q} \lor \lnot\metav{Q})) & \Leftrightarrow \metav{P}  && \text{(Simp)}\\
(\metav{P} \lor (\metav{Q} \land \lnot\metav{Q})) & \Leftrightarrow \metav{P}
\end{align*}
\begin{align*}
(\metav{P} \lor (\metav{Q} \lor \lnot\metav{Q})) & \Leftrightarrow (\metav{Q} \lor \lnot\metav{Q}) \\
(\metav{P} \land (\metav{Q} \land \lnot\metav{Q})) & \Leftrightarrow (\metav{Q} \land \lnot\metav{Q})
\end{align*}
\fi
